\chapter{时间序列、预测与状态空间基础}
\label{ch:time-series}

时间序列的难点不是把滞后变量放进回归，而是说明哪些信息在每个时点可用，以及哪些统计关系可以跨时间复用。本章从平稳性和创新开始，推导 AR(1) 的矩与预测误差，再把同一套问题带到 ARMA、单位根、GARCH 和状态空间模型。最后用伪回归、结构变化和随机切分的反例，说明预测结果为什么必须服从时间顺序。

读者可以把每个模型看成三件事的组合：状态如何演化、观测如何产生、验证如何向未来移动。只有三者一致，参数和预测区间才有可解释的含义。

\section{平稳性、创新与 AR(1) 预测}

严格平稳要求任意 $k$、任意时点与任意平移 $a$ 都有
\[
(X_{t_1},\ldots,X_{t_k})\overset d=(X_{t_1+a},\ldots,X_{t_k+a}).
\]
它约束全部有限维分布；弱平稳只约束常均值与仅依赖滞后的协方差。具有有限二阶矩的严格平稳过程必弱平稳，反向一般不成立。遍历性则连接时间平均与总体矩；平稳不自动保证一条长路径足以估计总体，例如 $X_t=Z$ 对所有 $t$ 恒定时严格平稳，却不能从单一路径恢复 $Z$ 的总体分布。

过程 $(X_t)$ 若满足 $\mathbb E[X_t]=\mu$、$\operatorname{Var}(X_t)<\infty$，且
\[
  \gamma(h)=\operatorname{Cov}(X_t,X_{t-h})
\]
只依赖滞后 $h$ 而不依赖 $t$，称为弱平稳；自相关函数记为 $\rho(h)=\gamma(h)/\gamma(0)$。白噪声 $(\varepsilon_t)$ 还满足零均值、常方差和不同期不相关；不相关不自动等于独立，金融收益中的波动聚集也不因线性自相关接近零而消失。

ARMA 模型把当前值写成过去状态和创新的有限线性组合：
\[
  X_t=c+\sum_{i=1}^{p}\phi_iX_{t-i}
      +\varepsilon_t+\sum_{j=1}^{q}\theta_j\varepsilon_{t-j}.
\]
令 $\mathcal F_t=\sigma(X_s:s\le t)$ 表示截至 $t$ 可用信息；若 $X_t$ 对 $\mathcal F_t$ 可测，则称过程适应于该信息流。创新 $\varepsilon_t$ 必须满足
\[
\varepsilon_t=X_t-\mathbb E(X_t\mid\mathcal F_{t-1}),\qquad
\mathbb E(\varepsilon_t\mid\mathcal F_{t-1})=0.
\]
因此它与任何平方可积的过去可测变量正交。拟合残差只是估计模型留下的样本差；只有模型与信息集声明成立时，才能把残差解释为创新。AR 根在单位圆外、MA 根在单位圆外分别对应常用的平稳与可逆条件。

这里必须区分层次：弱白噪声只要求零均值、常方差和跨期不相关，
弱白噪声不推出鞅差创新；创新还绑定指定信息集并要求
$\mathbb E(\varepsilon_t\mid\mathcal F_{t-1})=0$。独立同分布噪声更强，
但只有在其确实相对 $\mathcal F_{t-1}$ 独立时才自动给出这一条件均值。
ARMA 方程中的驱动噪声也不能仅凭“白噪声”标签便等同于观测滤过下的
一步预测创新；因果性、可逆性和所选滤过必须另行声明。

Wold 分解说明：纯非确定且协方差平稳的过程，可写成当前及过去创新的
无限移动平均。它解释了 ARMA 为何是有用近似，却不保证低阶 ARMA 必然
正确。白噪声只要求不相关；IID 噪声更强；高斯白噪声才因联合高斯而把
不相关提升为独立。

\MFQLead{AR(1) 的矩、相关与多步预测}

令 $X_t=\phi X_{t-1}+\varepsilon_t$，$|\phi|<1$，且创新方差为 $\sigma^2$。反复代入 $m$ 次得到
\[
X_t=\phi^mX_{t-m}+\sum_{j=0}^{m-1}\phi^j\varepsilon_{t-j}.
\]
平稳二阶矩与 $|\phi|<1$ 使第一项在 $L^2$ 中趋于零，故因果表示为 $X_t=\sum_{j\ge0}\phi^j\varepsilon_{t-j}$。创新跨期正交，于是几何级数给
\[
  \operatorname{Var}(X_t)=\frac{\sigma^2}{1-\phi^2},\qquad
  \gamma(h)=\phi^{|h|}\gamma(0).
\]
也可把模型乘以 $X_{t-h}$ 并取期望：当 $h\ge1$ 时，$\varepsilon_t$ 与 $X_{t-h}$ 正交，所以 Yule--Walker 递推为 $\gamma(h)=\phi\gamma(h-1)$；$h=0$ 时则有 $\gamma(0)=\phi^2\gamma(0)+\sigma^2$。

向前迭代 $h$ 步可把未来值拆成已知项与未来创新：
\[
X_{t+h}=\phi^hX_t+\sum_{j=0}^{h-1}\phi^j\varepsilon_{t+h-j}.
\]
给定 $\mathcal F_t$，和式条件均值为零、各项正交，因此
\[
  \widehat X_{t+h\mid t}=\phi^hX_t,
  \qquad
  \operatorname{Var}(X_{t+h}-\widehat X_{t+h\mid t})
  =\sigma^2\sum_{j=0}^{h-1}\phi^{2j}.
\]
当 $\phi=0.8,\sigma^2=1$ 时，理论方差为 $2.777778$；从 $X_t=1$ 出发的三步预测为 $0.512000$，预测误差方差为 $2.049600$。

\begin{example}[从预测误差到区间]
仍取 $\phi=0.8,\sigma^2=1$ 且 $X_t=1$。一步预测均值为 $0.8$、误差方差为 $1$；三步预测均值为 $0.512$、误差方差为
\[
  1+0.8^2+0.8^4=2.0496.
\]
若采用近似正态区间，三步预测的 95\% 区间为
$0.512\pm1.96\sqrt{2.0496}$。这个区间只反映已声明的创新分布与参数，不包含 $\phi$ 和 $\sigma^2$ 的估计不确定性；把估计参数当真值会低估预测风险。
\end{example}

\section{ARMA 的识别、估计与诊断}

对 $X_t-\phi X_{t-1}=\varepsilon_t+\theta\varepsilon_{t-1}$，AR 多项式 $1-\phi z$ 的根在单位圆外保证因果平稳表示，MA 多项式 $1+\theta z$ 的根在单位圆外保证创新可由观测历史唯一恢复。本例 $\phi=0.6,\theta=0.3$，根模分别为 $1.666667$ 与 $3.333333$。展开
\[
X_t=(1-\phi L)^{-1}(1+\theta L)\varepsilon_t
\]
得到 $\psi_0=1$、$\psi_j=(\phi+\theta)\phi^{j-1}$。创新正交使协方差成为系数乘积之和：
\[
\gamma(0)=\sigma^2\sum_{j\ge0}\psi_j^2
=\sigma^2\frac{1+\theta^2+2\phi\theta}{1-\phi^2}=2.265625,
\]
\[
\gamma(1)=\sigma^2\sum_{j\ge0}\psi_j\psi_{j+1}
=\sigma^2\frac{(\phi+\theta)(1+\phi\theta)}{1-\phi^2},
\qquad \rho(1)=\gamma(1)/\gamma(0)=0.732414.
\]
因此 ARMA 解析方差为 $2.265625$。
ACF 截尾与 PACF 截尾是理想总体识别线索，不是有限样本定阶定理。候选阶数还应比较残差、预测表现与参数稳定性；AIC 偏向预测效率，BIC 对模型维数惩罚更强，二者都不能修复数据泄漏。

\MFQLead{估计目标决定诊断含义}

条件平方和可以写成
\[
  Q(\phi,\theta)=\sum_{t=t_0}^T
  \widehat\varepsilon_t(\phi,\theta)^2,
\]
其中 $\widehat\varepsilon_t$ 由给定初值递推得到；高斯创新似然则为
\[
 \ell(\phi,\theta,\sigma^2)
 =-\frac12\sum_t\left[
 \log(2\pi\sigma^2)+
 \frac{\widehat\varepsilon_t^2}{\sigma^2}\right].
\]
优化时必须把根约束、方差正性和初值处理写进目标函数。AIC 的形式
$-2\ell+2k$、BIC 的形式 $-2\ell+k\log n$ 只是同一目标下的惩罚比较；若候选模型在不同窗口或不同损失上拟合，直接比较数值没有意义。Ljung--Box 统计量
\[
 Q_{\rm LB}=n(n+2)\sum_{h=1}^H
 \frac{\widehat\rho_h^2}{n-h}
\]
检验的是一组残差自相关，不能代替尾部、平方残差或样本外预测诊断。

\MFQLead{从估计到诊断：先声明目标函数}

对均值已处理的 AR$(p)$，条件最小二乘把 $X_t$ 回归于 $p$ 个滞后；若创新的条件矩正确，OLS 可一致，但近单位根会破坏常规有限样本近似。ARMA 的滞后创新不可见，常用条件平方和或高斯极大似然：给定参数递推创新与方差，再最小化平方和或最大化创新似然。GARCH 常用高斯准极大似然；即使尾部非高斯，在条件均值、条件方差和矩条件正确时仍可一致，但标准误应采用稳健版本。状态空间模型则直接用下文 Kalman 创新似然估计 $Q,R$ 与转移参数。

可复核流程是：只在训练窗口拟合候选模型，记录优化初值、参数约束、收敛状态与 Hessian/稳健标准误；随后才计算 AIC/BIC 和残差诊断，并把参数事先确定到下一验证窗口。近单位根、弱识别、局部最优、错误分布与结构变化都会让“优化器收敛”不等于“参数可解释”；Ljung--Box 的参考自由度还应扣除已估参数。估计不确定性最终要进入预测区间、参数漂移图和样本外损失，而不能只报一个点估计。

\section{单位根、差分与协整}

随机游走 $X_t=X_{t-1}+\varepsilon_t$ 的方差为 $t\sigma^2$，随时间增长；在时长 $1,10,100$ 上解析方差为 $1,10,100$。差分 $\Delta X_t=\varepsilon_t$ 平稳，但差分会丢失水平信息。单位根检验的低功效、确定性趋势和结构突变会改变结论，不能只根据一张 ACF 图决定差分次数。

\MFQLead{单位根检验不是“看起来不平稳”}

把 AR(1) 写成
\[
 \Delta X_t=\rho X_{t-1}+u_t,\qquad \rho=\phi-1.
\]
单位根假设是 $\rho=0$，但在该假设下 $X_{t-1}$ 的尺度随样本长度增长，普通 t 分布近似失效；ADF 检验通过加入 $\Delta X_{t-j}$ 控制短期相关，却仍需要带确定性项的专用临界值。KPSS 反过来把平稳性放在原假设，二者结论相反时应检查趋势、断点和检验功效，而不是挑一个更符合叙事的 p 值。

两个独立随机游走也可能在有限样本中产生高 $R^2$ 和显著斜率。原因不是存在稳定的长期关系，而是两个水平序列都携带累积噪声；差分后关系消失正是诊断线索。协整要求残差本身平稳，并且关系在样本外和经济机制上可解释；从大量配对中挑出最漂亮的残差后再做同一个检验，临界值和显著性都已经被选择污染。

若 $x_t,y_t$ 各自有单位根，但 $u_t=y_t-\beta x_t$ 平稳，则二者协整。设 $u_t=0.4u_{t-1}+e_t$、$\operatorname{Var}(e_t)=0.36$，则长期价差方差为 $0.428571$，且
\[
\Delta u_t=-0.6u_{t-1}+e_t,
\]
误差修正系数为 $-0.6$。协整是总体关系；先从大量配对中挑最平稳价差再原样检验，会产生选择偏差。经济机制、样本外稳定性和交易成本仍是必要检查。

固定种子价差实验使用长度 120 的滚动窗口逐次重估 AR 斜率，共得到 480 个估计；其最小值、中位数和最大值为 $(0.299681,0.459715,0.558824)$。这是滚动稳定性路径的摘要，不应与只在断点前后各拟合一次的分段估计混称。

\section{条件方差与状态空间}

收益可以近似无自相关而平方收益显著相关。GARCH$(1,1)$ 写为
\[
r_t=\sigma_tz_t,\qquad \sigma_t^2=\omega+\alpha r_{t-1}^2+\beta\sigma_{t-1}^2,
\]
其中 $z_t$ 相对 $\mathcal F_{t-1}$ 条件标准化：$\mathbb E(z_t\mid\mathcal F_{t-1})=0$、$\mathbb E(z_t^2\mid\mathcal F_{t-1})=1$，并令 $\sigma_t^2$ 为过去信息可测。因此 $\mathbb E(r_{t-1}^2)=\mathbb E(\sigma_{t-1}^2)$。若记平稳二阶矩 $v=\mathbb E\sigma_t^2$，对递推取期望便有
\[
v=\omega+(\alpha+\beta)v,\qquad
v=\frac{\omega}{1-\alpha-\beta}.
\]
当 $\omega>0$、$\alpha,\beta\ge0$ 且 $\alpha+\beta<1$ 时该矩有限。取 $(0.1,0.1,0.8)$，持续性为 $0.9$、无条件方差为 $1$，GARCH 半衰期为 $6.578813$。这不保证正态尾部或参数跨制度稳定；应检查标准化残差与其平方的剩余相关。

\MFQLead{状态空间模型把不可见状态与观测分开}

标量局部水平模型为
\[
  x_t=x_{t-1}+w_t,\quad w_t\sim\mathcal N(0,Q),
  \qquad y_t=x_t+v_t,\quad v_t\sim\mathcal N(0,R),
\]
并假设两类噪声跨期且彼此独立，也与初始状态独立。先验预测为 $x_{t\mid t-1}=x_{t-1\mid t-1}$、$P_{t\mid t-1}=P_{t-1\mid t-1}+Q$。设线性更新为 $x_{t\mid t}=x_{t\mid t-1}+K_t\nu_t$，其中 $\nu_t=y_t-x_{t\mid t-1}$。其条件误差方差为
\[
(1-K_t)^2P_{t\mid t-1}+K_t^2R.
\]
对 $K_t$ 求导并令其为零，得到最小均方误差增益；代回得到后验方差：
\[
  P_{t\mid t-1}=P_{t-1\mid t-1}+Q,
  \quad K_t=\frac{P_{t\mid t-1}}{P_{t\mid t-1}+R},
\]
\[
  x_{t\mid t}=x_{t\mid t-1}+K_t(y_t-x_{t\mid t-1}),
  \quad P_{t\mid t}=(1-K_t)P_{t\mid t-1}.
\]
取与全部噪声独立的高斯先验 $x_0\sim\mathcal N(0,1)$，即 $(x_{0\mid0},P_{0\mid0})=(0,1)$；令 $Q=0.25,R=1$，依次观测 $1,-0.5,0.25$，三次滤波均值依次为 $0.555556$、$0.084615$ 与 $0.152494$，末次增益为 $0.410431$。这些数由递推公式或独立参考实现给出，程序不得把滤波输出再当成自己的独立基准。

% Generated by tools/build_figures.py; edit figures/figure-specs.json instead.
\MFQFigure{tex/figures/generated/upper-ch12-state-space.pdf}{状态方程推进不可观测状态，观测方程只在当前时点产生带噪数据；两类不确定性不能混为一谈。}{upper-ch12-state-space}


每步预测条件于过去数据，所以 $\nu_t\mid\mathcal F_{t-1}\sim\mathcal N(0,S_t)$，其中 $S_t=P_{t\mid t-1}+R$。逐期条件密度相乘、取对数后得到高斯创新似然
\[
\ell=-\frac12\sum_t\{\log(2\pi)+\log S_t+\nu_t^2/S_t\}.
\]
本例创新对数似然为 $-4.260729$。
滤波只使用截至 $t$ 的数据；Rauch--Tung--Striebel 平滑使用全样本。在线性状态转移
$x_t=Fx_{t-1}+w_t$ 下，后向增益为
\[
 J_t=P_{t\mid t}F^T(P_{t+1\mid t})^{-1},
\]
平滑递推为
\[
 x_{t\mid T}=x_{t\mid t}+J_t(x_{t+1\mid T}-x_{t+1\mid t}),
\quad
 P_{t\mid T}=P_{t\mid t}+J_t(P_{t+1\mid T}-P_{t+1\mid t})J_t^T.
\]
三期平滑均值为 $(0.260771,0.128118,0.152494)$，所以不能把平滑状态当成实时特征。非高斯或非线性系统需要扩展/无迹 Kalman、粒子方法或其他近似，且似然比较必须基于同一观测协议。

\section{时间顺序、诊断与验证}

残差 ACF 和 Ljung--Box 检查一组滞后的线性相关，但拒绝并不说明该选哪种替代模型，不拒绝也不证明独立、同分布或无条件异方差。应同时检查残差平方、分布尾部、参数漂移和预测误差。定阶与调参只能在训练窗口内完成；外层滚动或扩展窗口负责评估，必要时在制度断点处报告分段结果。

\MFQLead{一个完整的时间顺序案例}

设有 1,000 个按日期排序的观测，先用前 600 个训练，接着 200 个验证，最后 200 个测试。AR 阶数和 GARCH 是否加入只能在前 800 个观测内决定；一旦选定，模型在前 800 个点上重估一次，再对最后 200 个点逐日生成一步预测。报告应同时给出训练残差诊断、验证损失、测试损失和预测区间覆盖率。若随机切分在训练中看到了测试制度，测试损失下降并不表示模型更好，而是评估边界已经被破坏。

\MFQLead{负例：时间顺序不是可随意打乱的行号}

两个独立随机游走在水平值上可能产生很高的 $R^2$，即伪回归；固定种子实验的水平回归 $R^2=0.344011$，差分回归 $R^2=0.000216$。高拟合度没有提供平稳残差、协整关系或经济机制。

再令 AR 系数在样本后段从 $0.85$ 变为 $-0.25$。断点前后分段 OLS 斜率为 $0.824923$ 与 $-0.215817$；随机打乱后切分又把新制度信息泄漏进训练集，随机切分 MSE 为 $1.436939$，时间切分 MSE 为 $2.226153$。更小的随机切分误差不是更可信的未来预测；研究报告应给出滚动或扩展窗口、断点诊断、训练与决策时点以及数据可得性。

\section{Kalman 更新为何是加权平均}

假设预测状态服从 $N(m,P)$，新观测为 $y=x+v$，$v\sim N(0,R)$ 且独立。后验密度与
\[
 \exp\left[-\tfrac12\left\{\frac{(x-m)^2}{P}+\frac{(y-x)^2}{R}\right\}\right]
\]
成比例。把关于 $x$ 的平方配方，后验精度为 $P^{-1}+R^{-1}$，后验均值为 $(R m+P y)/(P+R)=m+K(y-m)$，其中 $K=P/(P+R)$。这既重现前面的最小均方线性更新，也在高斯条件下证明它是完整条件期望。

$P=4,R=1,m=0,y=2$ 时，$K=0.8$，均值为 $1.6$、方差为 $0.8$。将测量噪声改为 $R=16$，增益降至 $0.2$，均值只移动到 $0.4$。notebook 在每一步显示预测均值、预测方差、创新与增益，并改变 $Q/R$。若只看一条平滑曲线，读者容易把“不信任观测”误认为“状态本来平稳”。

\noindent\textbf{延伸阅读。}ARMA 与预测参见\cite{brockwell2016timeseries}，状态空间与 Kalman 滤波参见\cite{harvey1990structural}。

\section{四级问题与即时练习}
\begingroup\small
\begin{enumerate}[nosep,leftmargin=*]
  \item \textbf{口述概念：}区分弱平稳、白噪声、创新、随机游走与状态空间中的隐状态。
  \item \textbf{笔试推导：}推导平稳 AR(1) 的方差、滞后协方差和 $h$ 步预测误差方差，并完成三次标量 Kalman 更新。
  \item \textbf{数值编程：}独立复现 AR(1) 矩、Kalman 手算值、随机游走水平/差分回归和结构变化切分误差。
  \item \textbf{研究判断：}一个随机切分表现优异的收益预测器上线后失效，应如何区分泄漏、结构变化与纯噪声？
\end{enumerate}
\endgroup

\subsection*{分级提示}
\begingroup\small
\begin{enumerate}[nosep,leftmargin=*]
  \item \textbf{口述概念：}先问均值和协方差是否依赖日历时间，再问误差相对于哪个信息集不可预测。
  \item \textbf{笔试推导：}对 AR(1) 反复代入并使用创新与过去不相关；Kalman 每步都先预测方差再计算增益。
  \item \textbf{数值编程：}解析矩和三步分数由程序外给出；模拟只检查容差，负例同时报告水平与差分。
  \item \textbf{研究判断：}冻结特征可得时点，使用滚动切分，按制度分段报告残差、自相关、断点与误差漂移。
\end{enumerate}
\endgroup


\noindent\textbf{本章要点。}平稳性规定哪些矩能跨时间复用，创新规定预测误差相对哪个信息集不可预测，状态空间递推把隐状态与测量噪声分开。任何时间序列结论都应附带时间切分、结构变化诊断、预测目标和失败边界；随机切分只能在明确不会泄漏时使用。
